\documentclass[10pt,a4paper]{amsart}
\usepackage[margin=19mm]{geometry}
\usepackage{amsmath,amssymb,mathtools}
\usepackage[scr=rsfs,cal=euler]{mathalfa}
\usepackage{xfrac}
\usepackage[unicode,hidelinks,bookmarks=false]{hyperref}
\newcommand{\ie}{i.e.\ }
\newcommand{\cf}{cf.\ }
\newcommand{\resp}{respectively\ }
\newcommand{\Subsection}{\subsection}
\providecommand{\nobreakdash}{\nobreak\hbox{-}}
\setlength{\emergencystretch}{2em}
\multlinegap0pt
\usepackage{mathtools}
\usepackage{bm}
\usepackage{mathrsfs}
\usepackage{textcomp}
\usepackage{manyfoot}
\usepackage{booktabs}
\usepackage{enumitem}
\newtheorem{theorem}{Theorem}
\newtheorem{corollary}{Corollary}
\newtheorem{lemma}{Lemma}
\newtheorem{proposition}{Proposition}
\newtheorem{example}{Example}
\newtheorem{definition}{Definition}
\newtheorem{problem}{Problem}
\newtheorem{remark}{Remark}
\begin{document}
\title[Cone Minimax Principles for Non-Selfadjoint Operator Pencils]{Cone Minimax Principles for Non-Selfadjoint Operator Pencils}
\author{Yavdat Il'yasov; Nur Valeev}
\date{}
\maketitle
\noindent\emph{Source and licence.} Cone Minimax Principles for Non-Selfadjoint Operator Pencils. Version arXiv:2606.31129v1 (CC BY 4.0). This material is distributed under the Creative Commons Attribution 4.0 International licence, http://creativecommons.org/licenses/by/4.0/, as recorded in the arXiv OAI licence metadata for this exact version and shown on the abstract page https://arxiv.org/abs/2606.31129v1. Full rights notice: licenses/arxiv-2606.31129-CC-BY-4.0.txt.\par\smallskip
\noindent\textit{Editorial scope.} Counted unit: the original \texttt{theorem} environment \texttt{thm:continuous-minimax-limit} at source lines 1139--1183 together with its complete proof at lines 1185--1361; the delivered document keeps source lines 468--1362 so that every referenced label and every citation resolves inside it. Native editable source is the paper's own concatenated TeX; the AMS wrapper is editorial formatting only. No publisher class, upstream script or unrelated artwork is redistributed.
\begin{lemma}[Minimax trapping]\label{lem:main}
	Every cone quasi-eigenvalue $\lambda_{*}$ satisfies
	\[
	\underline{\lambda} \leq \lambda_{*} \leq \overline{\lambda}.
	\]
	In particular, if the minimax condition holds, then $\hat\lambda$ is the
	unique cone quasi-eigenvalue.
\end{lemma}

\begin{proof}
	If $\lambda_{*}$ is realized by $\phi_{*},\psi_{*}\in\mathcal{S}\setminus\{0\}$,
	then $\overline{\lambda}(\phi_{*})=\lambda_{*}\leq\overline{\lambda}$ and
	$\underline{\lambda}\leq\underline{\lambda}(\psi_{*})=\lambda_{*}$.
\end{proof}

The next result shows that when the realizing vectors lie strictly inside the
cone, the trapping is sharp.

\begin{corollary}[Interior quasi-eigenvectors collapse the minimax levels]
	\label{cor:interior-eigenvalue-collapse}
	If $\lambda_{*}$ is a cone quasi-eigenvalue realized by
	$\phi_{*},\psi_{*}\in\mathcal{S}^{o}$, then
	$\underline{\lambda}=\lambda_{*}=\overline{\lambda}$.
	Consequently, failure of the minimax condition implies that no real eigenvalue
	admits a strictly positive right-left eigenpair.
\end{corollary}

\begin{proof}
	By Lemma~\ref{lem:main} it suffices to establish the reverse inequalities
	$\underline{\lambda}\geq\lambda_{*}$ and $\overline{\lambda}\leq\lambda_{*}$.
	We prove the first; the second is symmetric.
	
	Take any $\psi\in\mathcal{S}\setminus\{0\}$ and choose
	$\psi_{n}\in\mathcal{S}^{o}$ with $\psi_{n}\to\psi$ in $\mathcal{V}$.
	Since $\overline{\lambda}(\phi_{*})=\lambda_{*}$, we have
	$\mathcal{R}(\phi_{*},\psi_{n})\geq\lambda_{*}$ for all $n$. As
	$\phi_{*}\in\mathcal{S}^{o}$, condition $(P)$ gives
	$(\mathcal{G}\phi_{*},\psi)_{H}>0$, so the quotient is continuous in
	$\psi_{n}$ and the inequality survives in the limit:
	$\mathcal{R}(\phi_{*},\psi)\geq\lambda_{*}$. Since
	$\phi_{*}\in\mathcal{S}^{o}$, this means
	$\underline{\lambda}(\psi)\geq\mathcal{R}(\phi_{*},\psi)\geq\lambda_{*}$.
	As $\psi$ was arbitrary, $\underline{\lambda}\geq\lambda_{*}$.
	The proof of $\overline{\lambda}\le\lambda_{*}$ is the same, using
	$\psi_{*}\in\mathcal S^o$ and the identity
	$\underline{\lambda}(\psi_{*})=\lambda_{*}$.
\end{proof}

\begin{remark}[Strict positivity is essential]
	The conclusion fails for boundary eigenvectors. For
	$A=\operatorname{diag}(1,2)$, $G=I$, $\mathcal{S}=\mathbb{R}_{+}^{2}$, one
	finds $\underline{\lambda}=1<2=\overline{\lambda}$. The eigenvalue $2$ is
	realized by $e_{2}\in\partial\mathcal{S}$, yet the levels do not collapse,
	because $e_{2}$ is not an admissible strictly positive test direction.
\end{remark}

A complementary collapse mechanism operates even when the realizing vectors
reach the boundary of the cone, provided both principal levels are attained.

\begin{lemma}[Saddle collapse]\label{lem:saddle-collapse}
	Suppose $\overline{\lambda}$ and $\underline{\lambda}$ are finite and attained
	at $\phi^{*},\psi^{*}\in\mathcal{S}\setminus\{0\}$, with
	$(\mathcal{G}\phi^{*},\psi^{*})_{H}>0$. Then
	\[
	\overline{\lambda}
	\leq \mathcal{R}(\phi^{*},\psi^{*})
	\leq \underline{\lambda}.
	\]
	If, additionally, $\underline{\lambda}\leq\overline{\lambda}$ then
	\[
	\underline{\lambda}
	= \overline{\lambda}
	= \mathcal{R}(\phi^{*},\psi^{*}),
	\]
	and $(\phi^{*},\psi^{*})$ is a saddle point:
	\[
	\mathcal{R}(u,\psi^{*})
	\leq \mathcal{R}(\phi^{*},\psi^{*}) \leq
	\mathcal{R}(\phi^{*},v)
	\qquad \forall\,u,v\in\mathcal{S}^{o}.
	\]
\end{lemma}

\begin{proof}
	From $\overline{\lambda}(\phi^{*})=\overline{\lambda}$ we have
	$\mathcal{R}(\phi^{*},v)\geq\overline{\lambda}$ for all $v\in\mathcal{S}^{o}$.
	Choose $v_{n}\in\mathcal{S}^{o}$ with $v_{n}\to\psi^{*}$ in $\mathcal{V}$;
	the assumption $(\mathcal{G}\phi^{*},\psi^{*})_{H}>0$ allows passage to the
	limit, giving $\mathcal{R}(\phi^{*},\psi^{*})\geq\overline{\lambda}$.
	The symmetric argument using $\underline{\lambda}(\psi^{*})=\underline{\lambda}$
	yields $\mathcal{R}(\phi^{*},\psi^{*})\leq\underline{\lambda}$.
	If $\underline{\lambda}\leq\overline{\lambda}$, the three quantities are
	squeezed to equality, and the saddle inequalities follow directly from the
	definitions of $\overline{\lambda}(\phi^{*})$ and
	$\underline{\lambda}(\psi^{*})$. The additional ordering
	$\underline{\lambda}\leq\overline{\lambda}$ holds, for instance, whenever a
	cone quasi-eigenvalue exists, by Lemma~\ref{lem:main}.
\end{proof}

\begin{remark}[Relation to the spectrum]
	The cone minimax construction depends on the ordered triple
	$(\mathcal{L},\mathcal{G},\mathcal{S})$, not on $\mathcal{L}$ alone. When a
	strictly positive right-left eigenpair exists for a real eigenvalue
	$\lambda_{*}$, Corollary~\ref{cor:interior-eigenvalue-collapse} forces the
	minimax condition and identifies $\hat\lambda$ with $\lambda_{*}$.
	Contrapositively, failure of the minimax condition certifies the absence of
	any such eigenpair.
	
	When the relevant spectrum is non-real, as in the rotation example, the
	construction may still yield a real level $\hat\lambda$, but the corresponding
	quasi-pair is not a genuine spectral eigenpair; it is an intrinsically
	cone-theoretic object with no classical counterpart.
\end{remark}

%--------------------------------

\section{Discrete minimax formula}
\label{sec:approximation}

A notable feature of the finite-dimensional part is its generality. The
minimax identities below use only the ordered finite-dimensional cone
structure, strict positivity of the Galerkin weight, and a boundary-exclusion
condition. No ellipticity, cone invariance of the operator, resolvent
positivity, or compactness assumption is needed at this stage. These
additional analytic assumptions enter only later, when the discrete
minimax identities are connected with Galerkin limits of infinite-dimensional
operator pencils.

\medskip

\noindent\textbf{Elliptic realization.}\;
From this point on, when Galerkin approximation and elliptic examples are
considered, we use the following concrete realization of the abstract setting.
Let \(\Omega\subset\mathbb R^n\) be a bounded domain, set
\[
V:=H_0^1(\Omega),
\qquad
\mathcal V:=V^m,
\qquad
H:=[L^2(\Omega)]^m,
\]
and consider
\begin{equation}\label{eq:L_system}
	\mathcal L\mathbf u
	=
	-\sum_{i,j=1}^n \partial_j\!\bigl(A_{ij}(x)\partial_i\mathbf u\bigr)
	+
	\sum_{i=1}^n B_i(x)\partial_i\mathbf u
	+
	C(x)\mathbf u,
	\qquad
	\mathcal G\mathbf u=G(x)\mathbf u ,
\end{equation}
where
\[
A_{ij},\,B_i,\,C,\,G\in L^\infty(\Omega)^{m\times m}.
\]
No symmetry assumptions are imposed on these matrices. The operator
\(\mathcal L\) is understood through its weak realization
\(\mathcal L:\mathcal V\to\mathcal V^*\), given by
\[
\begin{aligned}
	\langle \mathcal L\mathbf u,\mathbf v\rangle
	&=
	\sum_{i,j=1}^n
	\int_\Omega
	\bigl(A_{ij}(x)\partial_i\mathbf u\bigr)\cdot
	\partial_j\mathbf v\,dx
	+
	\sum_{i=1}^n
	\int_\Omega
	\bigl(B_i(x)\partial_i\mathbf u\bigr)\cdot
	\mathbf v\,dx                                      \\
	&\quad
	+
	\int_\Omega
	\bigl(C(x)\mathbf u\bigr)\cdot \mathbf v\,dx,
	\qquad \mathbf u,\mathbf v\in\mathcal V .
\end{aligned}
\]
For admissible pairs we keep the notation
\begin{equation}\label{eq:Rel}
	\mathcal R(\mathbf u,\mathbf v)
	:=
	\frac{\langle \mathcal L\mathbf u,\mathbf v\rangle}
	{(\mathcal G\mathbf u,\mathbf v)_H}.
\end{equation}

We assume that $\Omega$ admits an interior polyhedral exhaustion: there exist
polyhedral subdomains $\Omega_r\subset\Omega$ with $\bigcup_{r\ge1}\Omega_r=\Omega$.
The index $r$ encodes both the choice of subdomain and the mesh size. On each
$\Omega_r$ we fix a shape-regular conforming $P_1$-triangulation $\mathcal{T}_r$
and let $V_r$ be the corresponding finite element space with zero boundary values
on $\partial\Omega_r$, extended by zero to $\Omega$, so that $V_r\subset V$. We
assume the standard Galerkin density property
\begin{equation}\label{eq:Wr-density}
	\forall\,\xi\in V \quad \exists\,\xi_r\in V_r \quad\text{such that}\quad
	\xi_r\longrightarrow\xi \quad\text{strongly in }V,
\end{equation}
which holds for standard interior polyhedral exhaustions and shape-regular
conforming spaces; see~\cite{ciaretRav,ciaret,Ern}.

Let $B_1,\ldots,B_{N_r}$ be the interior nodes of $\mathcal{T}_r$ with nodal
basis $\{\psi_1,\ldots,\psi_{N_r}\}$ satisfying $\psi_i(B_j)=\delta_{ij}$,
$\psi_i\ge0$, and $\psi_i=0$ on $\partial\Omega_r$. After extension by zero,
each $\psi_i$ belongs to $V\cap C(\overline\Omega)$.

\medskip

\noindent\textbf{Discrete cones and operators.}\;
Define the scalar nodal cones
\begin{align*}
	&S_r^o
	:= \Bigl\{\,u_r=\textstyle\sum_i u_i\psi_i\in V_r : u_i>0\Bigr\},
	\\
	&S_r
	:= \overline{S_r^o}^{\,V_r}
	= \Bigl\{\,u_r=\textstyle\sum_i u_i\psi_i\in V_r : u_i\ge0\Bigr\}.
\end{align*}
Thus $S_r^o$ is an open cone in $V_r$ and $S_r$ is a closed polyhedral cone
with basis $\{\psi_1,\ldots,\psi_{N_r}\}$. For systems we set
\[
\mathcal{V}_r := V_r^m, \qquad
\mathcal{S}_r^o := (S_r^o)^m, \qquad
\mathcal{S}_r := (S_r)^m,
\]
with basis vectors $\Phi_{i,k}:=\psi_i e_k$, $i=1,\ldots,N_r$,
$k=1,\ldots,m$. Since every element of $\mathcal{S}_r$ is continuous,
componentwise nonnegative, and belongs to $\mathcal{V}$, we have
$\mathcal{S}_r\subset\mathcal{S}$.

\medskip

For $\mathbf{u}_r=\sum_{i,k}u_k^i\Phi_{i,k}\in\mathcal{V}_r$, write
$\bar{\mathbf{u}}:=(u_k^i)_{i,k}\in\mathbb{R}^{mN_r}$ for its coefficient
vector, with all inequalities understood componentwise. Then
\[
\mathbf{u}_r\in\mathcal{S}_r \Longleftrightarrow \bar{\mathbf{u}}\ge0,
\qquad
\mathbf{u}_r\in\mathcal{S}_r^o \Longleftrightarrow \bar{\mathbf{u}}>0.
\]

For $\mathbf{u}\in\mathcal{C}^0$, the componentwise nodal interpolant is
\[
\mathcal{I}_r\mathbf{u}
:= \sum_{i=1}^{N_r}\sum_{k=1}^m u_k(B_i)\Phi_{i,k}.
\]
Since every node $B_i$ lies in $\Omega$, interpolation preserves cone membership:
\begin{equation}\label{eq:positive-interpolation}
	\mathcal{I}_r(\mathcal{S})\subset\mathcal{S}_r,
	\qquad
	\mathcal{I}_r(\mathcal{S}^o)\subset\mathcal{S}_r^o.
\end{equation}

Define the stiffness and weight matrices by
\[
(L_r)_{(j,\ell),(i,k)}
:= \langle\mathcal{L}\Phi_{i,k},\Phi_{j,\ell}\rangle,
\qquad
(G_r)_{(j,\ell),(i,k)}
:= (\mathcal{G}\Phi_{i,k},\Phi_{j,\ell})_H,
\]
so that $\bar{\mathbf{v}}^T L_r\bar{\mathbf{u}}=\langle\mathcal{L}\mathbf{u}_r,
\mathbf{v}_r\rangle$ and $\bar{\mathbf{v}}^T G_r\bar{\mathbf{u}}=
(\mathcal{G}\mathbf{u}_r,\mathbf{v}_r)_H$. Whenever the denominator is positive,
the discrete extended Rayleigh quotient
\begin{equation}\label{eq:discrete-Rayleigh}
	\mathcal{R}_r(\mathbf{u}_r,\mathbf{v}_r)
	:= \frac{\bar{\mathbf{v}}^T L_r\bar{\mathbf{u}}}
	{\bar{\mathbf{v}}^T G_r\bar{\mathbf{u}}}
\end{equation}
is the restriction of $\mathcal{R}$ to $\mathcal{V}_r\times\mathcal{V}_r$.

\subsection{Minimax characterization}

We work under the following two assumptions.

\begin{description}
	\item[$(P_r)$] \emph{Strict positivity of the Galerkin weight.}
	\[
	\bar{\mathbf{v}}^T G_r\bar{\mathbf{u}}>0
	\qquad
	\forall\,(\mathbf{u}_r,\mathbf{v}_r)\in\mathcal{D}_r,
	\]
	where $\mathcal{D}_r:=(\mathcal{S}_r\setminus\{\mathbf{0}\})\times\mathcal{S}_r^o
	\cup\,\mathcal{S}_r^o\times(\mathcal{S}_r\setminus\{\mathbf{0}\})$.
	
	\item[$(B_r)$] \emph{Boundary exclusion.} For $M=L_r,L_r^T$,
	\[
	\bar{\mathbf{w}}\ge0,\quad
	\bar{\mathbf{w}}\ne0,\quad
	M\bar{\mathbf{w}}\ge0
	\quad\Longrightarrow\quad
	\bar{\mathbf{w}}>0.
	\]
\end{description}

Define the four mixed minimax levels
\begin{align}
	\overline\lambda_r
	&:= \sup_{\mathbf{u}\in\mathcal{S}_r\setminus\{\mathbf{0}\}}
	\inf_{\mathbf{v}\in\mathcal{S}_r^o}
	\mathcal{R}_r(\mathbf{u},\mathbf{v}),
	&
	\overline\Lambda_r
	&:= \inf_{\mathbf{v}\in\mathcal{S}_r^o}
	\sup_{\mathbf{u}\in\mathcal{S}_r\setminus\{\mathbf{0}\}}
	\mathcal{R}_r(\mathbf{u},\mathbf{v}),
	\label{eq:discrete-upper-minimax}\\
	\underline\Lambda_r
	&:= \sup_{\mathbf{u}\in\mathcal{S}_r^o}
	\inf_{\mathbf{v}\in\mathcal{S}_r\setminus\{\mathbf{0}\}}
	\mathcal{R}_r(\mathbf{u},\mathbf{v}),
	&
	\underline\lambda_r
	&:= \inf_{\mathbf{v}\in\mathcal{S}_r\setminus\{\mathbf{0}\}}
	\sup_{\mathbf{u}\in\mathcal{S}_r^o}
	\mathcal{R}_r(\mathbf{u},\mathbf{v}).
	\label{eq:discrete-lower-minimax}
\end{align}

\begin{lemma}[Discrete minimax principle]\label{lem:discrete-minimax}
	Assume $(P_r)$. Then both pairs of levels coincide and are finite:
	\[
	-\infty < \overline\lambda_r = \overline\Lambda_r < +\infty,
	\qquad
	-\infty < \underline\lambda_r = \underline\Lambda_r < +\infty.
	\]
	Moreover, there exist $\mathbf{u}_r,\mathbf{v}_r\in\mathcal{S}_r\setminus\{\mathbf{0}\}$
	such that
	\begin{equation}\label{eq:minmaxD2}
		\overline\lambda_r
		= \inf_{\mathbf{v}\in\mathcal{S}_r^o}\mathcal{R}_r(\mathbf{u}_r,\mathbf{v}),
		\qquad
		\underline\lambda_r
		= \sup_{\mathbf{u}\in\mathcal{S}_r^o}\mathcal{R}_r(\mathbf{u},\mathbf{v}_r).
	\end{equation}
\end{lemma}

\begin{proof}
	By homogeneity, all variables may be normalized to the compact convex set
	\[
	K_r := \Bigl\{\mathbf{u}\in\mathcal{S}_r :
	\textstyle\sum_{i,k}u_k^i=1\Bigr\},
	\qquad K_r^o := K_r\cap\mathcal{S}_r^o,
	\]
	where $K_r^o$ is convex and dense in $K_r$. By $(P_r)$, the denominator of
	$\mathcal{R}_r$ is positive on $K_r\times K_r^o$ and $K_r^o\times K_r$.
	
	The quotient is continuous and linear-fractional in each variable;
	hence it is quasiconcave in its first variable and quasiconvex in its
	second. Sion's minimax theorem~\cite{Sion1958} therefore gives
	\[
	\overline\lambda_r=\overline\Lambda_r.
	\]
	Applying it to
	\((\mathbf v,\mathbf u)\mapsto-\mathcal R_r(\mathbf u,\mathbf v)\)
	gives
	\[
	\underline\lambda_r=\underline\Lambda_r.
	\]
	Finiteness follows because, for any fixed
	$\mathbf{u}^0,\mathbf{v}^0\in K_r^o$, the denominators
	$\bar{\mathbf{v}}^T G_r\bar{\mathbf{u}}^0$ and
	$(\bar{\mathbf{v}}^0)^T G_r\bar{\mathbf{u}}$ are bounded away from zero on
	$K_r$. Finally, the maps
	$\mathbf{u}\mapsto\inf_{\mathbf{v}\in K_r^o}\mathcal{R}_r(\mathbf{u},\mathbf{v})$
	and $\mathbf{v}\mapsto\sup_{\mathbf{u}\in K_r^o}\mathcal{R}_r(\mathbf{u},\mathbf{v})$
	are upper and lower semicontinuous on the compact set $K_r$, respectively, so
	their extrema are attained, proving~\eqref{eq:minmaxD2}.
\end{proof}

\begin{lemma}[Discrete minimax collapse]\label{lem:discrete-collapse}
	Assume $(P_r)$ and $(B_r)$, and let $\mathbf{u}_r,\mathbf{v}_r$ be as in
	Lemma~\ref{lem:discrete-minimax}. If $\overline\lambda_r\ge0$, then
	$\mathbf{u}_r,\mathbf{v}_r\in\mathcal{S}_r^o$,
	\[
	\hat\lambda_r := \overline\lambda_r = \underline\lambda_r,
	\qquad
	\mathcal{R}_r(\mathbf{u}_r,\mathbf{v}_r) = \hat\lambda_r,
	\]
	and the pair $(\mathbf{u}_r,\mathbf{v}_r)$ satisfies the discrete eigenvalue
	equations
	\begin{equation}\label{eq:uravn}
		L_r\bar{\mathbf{u}}_r = \hat\lambda_r G_r\bar{\mathbf{u}}_r,
		\qquad
		L_r^T\bar{\mathbf{v}}_r = \hat\lambda_r G_r^T\bar{\mathbf{v}}_r.
	\end{equation}
\end{lemma}

\begin{proof}
	We first prove that $\mathbf{u}_r$ is strictly positive. By
	\eqref{eq:minmaxD2},
	$\mathcal{R}_r(\mathbf{u}_r,\mathbf{v})\ge\overline\lambda_r$
	for all $\mathbf{v}\in\mathcal{S}_r^o$. Since the denominator is positive,
	\[
	\bar{\mathbf v}^{T}(L_r-\overline\lambda_r G_r)\bar{\mathbf u}_r\ge0
	\qquad \forall\,\mathbf v\in\mathcal S_r^o,
	\]
	and hence $(L_r-\overline\lambda_r G_r)\bar{\mathbf{u}}_r\ge0$ componentwise.
	Also, $(P_r)$ gives $G_r\bar{\mathbf{u}}_r\ge0$: indeed,
	$\bar{\mathbf v}^{T}G_r\bar{\mathbf u}_r>0$ for all
	$\mathbf v\in\mathcal S_r^o$, and passage to boundary directions yields
	componentwise nonnegativity. Since $\overline\lambda_r\ge0$, it follows that
	$L_r\bar{\mathbf{u}}_r\ge0$. Therefore $(B_r)$ gives
	$\bar{\mathbf{u}}_r>0$, that is, $\mathbf{u}_r\in\mathcal{S}_r^o$.
	
	Since $\mathbf{u}_r\in\mathcal{S}_r^o$, $(P_r)$ keeps the denominator positive
	for $\mathbf v\in\mathcal S_r\setminus\{\mathbf0\}$, and the infimum in
	\eqref{eq:minmaxD2} extends by continuity from $\mathcal{S}_r^o$ to
	$\mathcal{S}_r\setminus\{\mathbf{0}\}$. Thus
	$\underline\Lambda_r\ge\overline\lambda_r$. Conversely, for every
	$\mathbf u\in\mathcal{S}_r^o$,
	\[
	\inf_{\mathbf v\in\mathcal{S}_r\setminus\{\mathbf0\}}
	\mathcal R_r(\mathbf u,\mathbf v)
	\le
	\inf_{\mathbf v\in\mathcal{S}_r^o}
	\mathcal R_r(\mathbf u,\mathbf v)
	\le
	\overline\lambda_r,
	\]
	and therefore $\underline\Lambda_r\le\overline\lambda_r$. Hence
	\[
	\underline\lambda_r=\underline\Lambda_r=\overline\lambda_r=\overline\Lambda_r=:\hat\lambda_r.
	\]
	
	It remains to identify the saddle value and prove strict positivity of
	$\mathbf{v}_r$. From \eqref{eq:minmaxD2},
	$\mathcal{R}_r(\mathbf{u},\mathbf{v}_r)\le\hat\lambda_r$ for all
	$\mathbf{u}\in\mathcal{S}_r^o$, while the extended infimum gives
	$\mathcal{R}_r(\mathbf{u}_r,\mathbf{v}_r)\ge\hat\lambda_r$. Since
	$\mathbf{u}_r\in\mathcal{S}_r^o$, the first inequality applies to
	$\mathbf{u}_r$ as well, and therefore
	\[
	\mathcal{R}_r(\mathbf{u}_r,\mathbf{v}_r)=\hat\lambda_r.
	\]
	
	Now $\mathcal{R}_r(\mathbf{u},\mathbf{v}_r)\le\hat\lambda_r$ for all
	$\mathbf u\in\mathcal S_r^o$, with equality at the strictly positive point
	$\mathbf u_r$. Hence, for every coefficient direction $\boldsymbol\xi$ and
	all sufficiently small \(t\), the vector
	\(\bar{\mathbf u}_r+t\boldsymbol\xi\) remains positive, and differentiation
	at \(t=0\) gives
	\[
	L_r^T\bar{\mathbf{v}}_r=\hat\lambda_r G_r^T\bar{\mathbf{v}}_r.
	\]
	By $(P_r)$, $G_r^T\bar{\mathbf{v}}_r\ge0$: indeed,
	$\bar{\mathbf v}_r^T G_r\bar{\mathbf u}>0$ for all
	$\mathbf u\in\mathcal S_r^o$, and passage to boundary directions yields
	componentwise nonnegativity. Since $\hat\lambda_r\ge0$, this implies
	$L_r^T\bar{\mathbf{v}}_r\ge0$. Applying $(B_r)$ to $L_r^T$ yields
	$\bar{\mathbf{v}}_r>0$, that is, $\mathbf{v}_r\in\mathcal{S}_r^o$.
	
	Finally, since $\mathbf{v}_r$ is also strictly positive and
	$\mathcal{R}_r(\mathbf{u}_r,\mathbf{v})\ge\hat\lambda_r$ for all
	$\mathbf v\in\mathcal S_r^o$, with equality at $\mathbf v_r$, differentiating
	with respect to arbitrary coefficient directions gives
	\[
	L_r\bar{\mathbf{u}}_r=\hat\lambda_r G_r\bar{\mathbf{u}}_r.
	\]
	This proves \eqref{eq:uravn}.
\end{proof}


\section{Existence of a left cone quasi-eigenvector}
\label{sec:left-minimax-attainment}

We first treat the lower/left level separately, because this part of the
construction admits a compactness argument under weaker assumptions and will
be used later in identifying the Galerkin limit. Throughout this section we
use the standing compactness assumptions that $\mathcal V$ is reflexive, the
embedding $\mathcal V\hookrightarrow H$ is compact, and $\mathcal S$ is weakly
closed in $\mathcal V$.

We assume the uniform ellipticity condition: there exists $\theta>0$ such that,
for a.e.\ $x\in\Omega$,
\begin{equation}\label{eq:unifellipt}
	\sum_{i,j=1}^n
	\langle A_{ij}(x)\xi_i,\xi_j\rangle_{\mathbb{R}^m}
	\ge \theta|\xi|^2
	\qquad
	\forall\,\xi=(\xi_1,\ldots,\xi_n)\in(\mathbb{R}^m)^n.
\end{equation}
This implies the G{\aa}rding inequality: there exist $c>0$ and $C\ge0$ such that
\[
\langle\mathcal{L}\mathbf{u},\mathbf{u}\rangle
\ge c\|\mathbf{u}\|_{\mathcal{V}}^2 - C\|\mathbf{u}\|_H^2
\qquad
\forall\,\mathbf{u}\in\mathcal{V}.
\]
Recall that for $\mathbf{v}\in\mathcal{S}\setminus\{\mathbf{0}\}$ the lower
one-sided level and the lower principal level are
\[
\underline\lambda(\mathbf{v})
:= \sup_{\mathbf{u}\in\mathcal{S}^o}\mathcal{R}(\mathbf{u},\mathbf{v}),
\qquad
\underline\lambda
:= \inf_{\mathbf{v}\in\mathcal{S}\setminus\{\mathbf{0}\}}\underline\lambda(\mathbf{v}).
\]

\begin{lemma}[Compactness of normalized left sublevels]
	\label{lem:left-sublevel-compactness}
	Assume $\mathrm{(P)}$ and \eqref{eq:unifellipt}. Let
	$(\mathbf{v}_j)\subset\mathcal{S}$ and $(\mu_j)\subset\mathbb{R}$ satisfy
	\[
	\|\mathbf{v}_j\|_H=1,
	\qquad \sup_j\mu_j<+\infty,
	\qquad
	\langle\mathcal{L}\mathbf{v}_j,\mathbf{v}_j\rangle
	\le \mu_j(\mathcal{G}\mathbf{v}_j,\mathbf{v}_j)_H.
	\]
	Then $(\mathbf{v}_j)$ is bounded in $\mathcal{V}$, and up to a subsequence
	\[
	\mathbf{v}_j\rightharpoonup\mathbf{v}
	\quad\text{weakly in }\mathcal{V},
	\qquad
	\mathbf{v}_j\to\mathbf{v}
	\quad\text{strongly in }H,
	\]
	for some $\mathbf{v}\in\mathcal{S}$ with $\|\mathbf{v}\|_H=1$.
\end{lemma}

\begin{proof}
	Let $C_G>0$ be such that
	\[
	|(\mathcal G\mathbf u,\mathbf w)_H|
	\le C_G\|\mathbf u\|_H\|\mathbf w\|_H
	\qquad \forall\,\mathbf u,\mathbf w\in H.
	\]
	By the density of $\mathcal{S}^o$ in $\mathcal{S}\setminus\{\mathbf{0}\}$,
	for each $j$ we may choose $\mathbf w_{j,n}\in\mathcal S^o$ such that
	$\mathbf w_{j,n}\to\mathbf v_j$ in $\mathcal V$. Since
	$\mathbf v_j\in\mathcal S\setminus\{\mathbf0\}$ and $\mathrm{(P)}$ gives
	$(\mathcal G\mathbf v_j,\mathbf w_{j,n})_H>0$, passage to the limit yields
	$(\mathcal G\mathbf v_j,\mathbf v_j)_H\ge0$. Together with the sublevel
	inequality and $\|\mathbf{v}_j\|_H=1$, this gives
	\[
	\langle\mathcal{L}\mathbf{v}_j,\mathbf{v}_j\rangle
	\le \max\bigl\{\sup_j\mu_j,0\bigr\}\,C_G,
	\]
	so the G{\aa}rding inequality bounds $(\mathbf{v}_j)$ in $\mathcal{V}$.
	The remaining assertions follow from the reflexivity of $\mathcal{V}$, the
	compact embedding $\mathcal{V}\hookrightarrow H$, and the weak closedness of
	$\mathcal{S}$.
\end{proof}

\begin{theorem}[Attainment of the lower principal level]
	\label{thm:general-left-attainment}
	Assume $\mathrm{(P)}$ and \eqref{eq:unifellipt}. If $\underline\lambda<+\infty$,
	then $\underline\lambda\in\mathbb{R}$ and there exists
	$\mathbf{v}^*\in\mathcal{S}\setminus\{\mathbf{0}\}$ with $\|\mathbf{v}^*\|_H=1$
	such that $\underline\lambda(\mathbf{v}^*)=\underline\lambda$ and
	\begin{equation}\label{eq:continuous-left-quasi}
		\langle\mathcal{L}\mathbf{u},\mathbf{v}^*\rangle
		\le \underline\lambda\,(\mathcal{G}\mathbf{u},\mathbf{v}^*)_H
		\qquad \forall\,\mathbf{u}\in\mathcal{S}.
	\end{equation}
\end{theorem}

\begin{proof}
	By homogeneity, there exists a minimizing sequence
	$(\mathbf{v}_k)\subset\mathcal{S}$ with $\|\mathbf{v}_k\|_H=1$ and finite
	values $\lambda_k:=\underline\lambda(\mathbf{v}_k)$ such that
	$\lambda_k\to\underline\lambda$. The sequence $(\lambda_k)$ is bounded above.
	For each fixed $k$, the definition of $\lambda_k$ gives
	\[
	\langle\mathcal{L}\mathbf{u},\mathbf{v}_k\rangle
	\le \lambda_k(\mathcal{G}\mathbf{u},\mathbf{v}_k)_H
	\qquad \forall\,\mathbf{u}\in\mathcal{S}^o.
	\]
	By density and continuity, this inequality extends to all
	$\mathbf{u}\in\mathcal{S}$. Taking $\mathbf{u}=\mathbf{v}_k$ and applying
	Lemma~\ref{lem:left-sublevel-compactness}, we extract a subsequence with
	$\mathbf{v}_k\rightharpoonup\mathbf{v}^*$ weakly in $\mathcal{V}$ and
	$\mathbf{v}_k\to\mathbf{v}^*$ strongly in $H$, where
	$\mathbf{v}^*\in\mathcal{S}$ and $\|\mathbf{v}^*\|_H=1$.
	
	If $\underline\lambda=-\infty$, then for any fixed $\mathbf{u}_0\in\mathcal{S}^o$,
	\[
	\lambda_k \ge \mathcal{R}(\mathbf{u}_0,\mathbf{v}_k)
	\longrightarrow \mathcal{R}(\mathbf{u}_0,\mathbf{v}^*)\in\mathbb{R},
	\]
	a contradiction. Hence $\underline\lambda\in\mathbb{R}$.
	
	Passing to the limit in the sublevel inequality gives
	\[
	\langle\mathcal{L}\mathbf{u},\mathbf{v}^*\rangle
	\le \underline\lambda(\mathcal{G}\mathbf{u},\mathbf{v}^*)_H
	\qquad \forall\,\mathbf{u}\in\mathcal{S}^o.
	\]
	Dividing by the positive denominator gives
	$\underline\lambda(\mathbf{v}^*)\le\underline\lambda$. The reverse inequality
	holds by definition, so $\underline\lambda(\mathbf{v}^*)=\underline\lambda$.
	Density and continuity then extend the inequality to all
	$\mathbf{u}\in\mathcal{S}$, proving \eqref{eq:continuous-left-quasi}.
\end{proof}



\section{Continuous minimax limit}
\label{sec:pde}

Under the discrete positivity and boundary-exclusion assumptions, the
finite-dimensional theory yields exact positive right-left eigenpairs and
exact minimax identities on each discrete cone. We now pass to the Galerkin
limit. The proof shows that every convergent subsequence of discrete minimax
values produces a nontrivial right-left eigenpair of the continuous pencil;
strong positivity is then used to identify this limit with the continuous cone
minimax value.

Throughout this section we assume that \(\mathcal V\) is reflexive, the
embedding \(\mathcal V\hookrightarrow H\) is compact, and \(\mathcal S\) is
weakly closed in \(\mathcal V\). Introduce the dual cone
\[
\mathcal S^*
:=
\{\mathbf F\in\mathcal V^*:
\langle\mathbf F,\mathbf w\rangle\ge0
\ \forall\,\mathbf w\in\mathcal S\},
\]
and assume:

\begin{description}
	\item[$\mathrm{(SP)}$] \emph{Strong positivity of \(\mathcal L\).}
	If \(\mathbf w\in\mathcal S\setminus\{\mathbf0\}\) and
	\(\mathcal L\mathbf w\in\mathcal S^*\), then
	\(\mathbf w\in\mathcal S^o\).
	
	\item[$\mathrm{(SP^*)}$] \emph{Strong positivity of \(\mathcal L^*\).}
	If \(\mathbf w\in\mathcal S\setminus\{\mathbf0\}\) and
	\(\mathcal L^*\mathbf w\in\mathcal S^*\), then
	\(\mathbf w\in\mathcal S^o\).
\end{description}

\begin{remark}[Relation with the maximum principle]
	Conditions \(\mathrm{(SP)}\) and \(\mathrm{(SP^*)}\) are cone formulations
	of the strong positivity consequence of the maximum principle. For the
	standard positive cone, the inclusion
	\(\mathcal L\mathbf w\in\mathcal S^*\) means that
	\(\mathcal L\mathbf w\ge0\) in the weak sense. Thus \(\mathrm{(SP)}\) says
	that every nontrivial nonnegative supersolution is strictly positive:
	\[
	\mathbf w\in\mathcal S\setminus\{\mathbf0\},
	\qquad
	\mathcal L\mathbf w\ge0
	\quad\Longrightarrow\quad
	\mathbf w\in\mathcal S^o .
	\]
	In the scalar elliptic case this is precisely the usual strong maximum
	principle; see, for example, \cite{ProtterWeinberger}. For systems it is the
	corresponding strong positivity property with respect to the chosen cone,
	and it typically reflects an irreducible coupling of the components. The
	adjoint condition \(\mathrm{(SP^*)}\) is the same requirement for the left
	problem and is needed to obtain strict positivity of the limiting left
	eigenvector.
	
	The present formulation is also consistent with the classical
	maximum-principle approach to principal eigenvalues. In particular,
	Donsker--Varadhan \cite{Donsker} obtained a variational formula for the
	principal eigenvalue of a nonsymmetric operator satisfying a maximum
	principle, and the relation between principal eigenvalues and maximum
	principles for elliptic operators was further developed by
	Berestycki--Nirenberg--Varadhan \cite{BerestNienbVardan}. In the present
	argument the maximum-principle input has a more specific role: it is not used
	to define the principal value or to reduce the pencil to a positive
	resolvent. Rather, after the Galerkin minimax limit has produced nonnegative
	right-left quasi-eigenvectors, \(\mathrm{(SP)}\) and \(\mathrm{(SP^*)}\) are
	used to upgrade them to strictly positive right-left eigenvectors.
\end{remark}


\begin{theorem}[Convergence of the finite-dimensional minimax identities]
	\label{thm:continuous-minimax-limit}
	Assume \(\mathrm{(P)}\), \eqref{eq:unifellipt}, and the standing compactness
	hypotheses of this section. Suppose that, for every \(r\), the discrete
	minimax collapse lemma applies, and assume that, for some \(C>0\),
	\begin{equation}\label{eq:estM}
		0<\hat\lambda_r\le C
		\qquad
		\forall\,r\ge1.
	\end{equation}
	If, in addition, \(\mathrm{(SP)}\) and \(\mathrm{(SP^*)}\) hold, then the pencil
	\((\mathcal L,\mathcal G)\) satisfies the minimax condition
	\begin{equation}\label{eq:mimaPrin1}
		\hat\lambda
		:=
		\sup_{\mathbf u\in\mathcal S\setminus\{\mathbf0\}}
		\inf_{\mathbf v\in\mathcal S^o}\mathcal R(\mathbf u,\mathbf v)
		=
		\inf_{\mathbf v\in\mathcal S\setminus\{\mathbf0\}}
		\sup_{\mathbf u\in\mathcal S^o}\mathcal R(\mathbf u,\mathbf v)
		\ge
		0,
	\end{equation}
	and \(\hat\lambda_r\to\hat\lambda\). Let
	\(\mathbf u_r,\mathbf v_r\in\mathcal S_r^o\) be the corresponding discrete
	right-left eigenpairs, normalized by
	\[
	\|\mathbf u_r\|_{\mathcal V}
	=
	\|\mathbf v_r\|_{\mathcal V}
	=
	1.
	\]
	Then these normalized pairs have weak-in-\(\mathcal V\) and strong-in-\(H\)
	cluster points, and every such cluster point \((\mathbf u_*,\mathbf v_*)\)
	belongs to \(\mathcal S^o\times\mathcal S^o\), solves
	\[
	(\mathcal L-\hat\lambda\mathcal G)\mathbf u_*=0,
	\qquad
	(\mathcal L-\hat\lambda\mathcal G)^*\mathbf v_*=0,
	\]
	and satisfies \(\mathcal R(\mathbf u_*,\mathbf v_*)=\hat\lambda\). Hence
	every such cluster point is a principal right-left quasi-pair realizing
	the continuous minimax identity.
\end{theorem}

\begin{proof}
	By the discrete minimax collapse lemma,
	\(\mathbf u_r,\mathbf v_r\in\mathcal S_r^o\) satisfy
	\begin{equation}\label{eq:discrete-minimax-eigenpairs}
		L_r\bar{\mathbf u}_r
		=
		\hat\lambda_r G_r\bar{\mathbf u}_r,
		\qquad
		L_r^T\bar{\mathbf v}_r
		=
		\hat\lambda_r G_r^T\bar{\mathbf v}_r.
	\end{equation}
	We first show that the \(H\)-limits of normalized subsequences cannot vanish.
	If \(\|\mathbf u_r\|_H\to0\) along a subsequence, then testing the first
	equation in \eqref{eq:discrete-minimax-eigenpairs} by \(\mathbf u_r\) gives
	\[
	\langle\mathcal L\mathbf u_r,\mathbf u_r\rangle
	=
	\hat\lambda_r(\mathcal G\mathbf u_r,\mathbf u_r)_H
	\to
	0,
	\]
	which contradicts the G{\aa}rding inequality and the normalization
	\(\|\mathbf u_r\|_{\mathcal V}=1\). The same argument, using
	\[
	\bar{\mathbf v}_r^T L_r^T\bar{\mathbf v}_r
	=
	\bar{\mathbf v}_r^T L_r\bar{\mathbf v}_r,
	\]
	excludes \(\|\mathbf v_r\|_H\to0\).
	
	Let \(\lambda_*\) be any accumulation point of \((\hat\lambda_r)\). Passing
	to a subsequence, we may assume that \(\hat\lambda_r\to\lambda_*\). By
	\eqref{eq:estM}, \(\lambda_*\ge0\). By reflexivity, the compact embedding
	\(\mathcal V\hookrightarrow H\), and the normalization, we may further
	assume that
	\[
	\mathbf u_r\rightharpoonup\mathbf u_*,
	\qquad
	\mathbf v_r\rightharpoonup\mathbf v_*
	\quad\text{weakly in }\mathcal V,
	\]
	and
	\[
	\mathbf u_r\to\mathbf u_*,
	\qquad
	\mathbf v_r\to\mathbf v_*
	\quad\text{strongly in }H.
	\]
	The weak closedness of \(\mathcal S\), together with the nonvanishing of the
	\(H\)-limits, gives
	\[
	\mathbf u_*,\mathbf v_*
	\in
	\mathcal S\setminus\{\mathbf0\}.
	\]
	
	Let \(\mathbf z\in\mathcal V\). Choose
	\(\mathbf z_r\in\mathcal V_r\) as in \eqref{eq:Wr-density}. Passing to the
	limit in
	\[
	\langle\mathcal L\mathbf u_r,\mathbf z_r\rangle
	=
	\hat\lambda_r(\mathcal G\mathbf u_r,\mathbf z_r)_H
	\]
	gives
	\[
	\langle\mathcal L\mathbf u_*,\mathbf z\rangle
	=
	\lambda_*(\mathcal G\mathbf u_*,\mathbf z)_H.
	\]
	Similarly, the second equation in \eqref{eq:discrete-minimax-eigenpairs}
	gives
	\[
	\langle\mathcal L\mathbf z_r,\mathbf v_r\rangle
	=
	\hat\lambda_r(\mathcal G\mathbf z_r,\mathbf v_r)_H,
	\]
	and passage to the limit yields
	\[
	\langle\mathcal L\mathbf z,\mathbf v_*\rangle
	=
	\lambda_*(\mathcal G\mathbf z,\mathbf v_*)_H.
	\]
	Hence
	\begin{equation}\label{eq:LimEq}
		(\mathcal L-\lambda_*\mathcal G)\mathbf u_*=0,
		\qquad
		(\mathcal L-\lambda_*\mathcal G)^*\mathbf v_*=0.
	\end{equation}
	
	Condition \(\mathrm{(P)}\) and the density of
	\(\mathcal S^o\) in \(\mathcal S\setminus\{\mathbf0\}\) imply
	\[
	(\mathcal G\mathbf u_*,\mathbf w)_H\ge0,
	\qquad
	(\mathcal G\mathbf w,\mathbf v_*)_H\ge0
	\qquad
	\forall\,\mathbf w\in\mathcal S.
	\]
	Equivalently,
	\[
	\mathcal G\mathbf u_*,
	\mathcal G^*\mathbf v_*
	\in
	\mathcal S^*.
	\]
	Since \(\lambda_*\ge0\), \eqref{eq:LimEq} gives
	\[
	\mathcal L\mathbf u_*,
	\mathcal L^*\mathbf v_*
	\in
	\mathcal S^*.
	\]
	Thus \(\mathrm{(SP)}\) and \(\mathrm{(SP^*)}\) yield
	\[
	\mathbf u_*,\mathbf v_*\in\mathcal S^o.
	\]
	
	The two eigenvalue equations imply
	\[
	\mathcal R(\mathbf u_*,\mathbf v)=\lambda_*
	\quad
	\forall\,\mathbf v\in\mathcal S^o,
	\qquad
	\mathcal R(\mathbf u,\mathbf v_*)=\lambda_*
	\quad
	\forall\,\mathbf u\in\mathcal S^o.
	\]
	Therefore
	\begin{equation}\label{eq:limit-quasi-pair}
		\overline\lambda(\mathbf u_*)
		=
		\lambda_*
		=
		\underline\lambda(\mathbf v_*),
		\qquad
		\underline\lambda\le\lambda_*\le\overline\lambda,
	\end{equation}
	by Lemma~\ref{lem:main}. Since
	\(\mathbf u_*,\mathbf v_*\in\mathcal S^o\),
	Corollary~\ref{cor:interior-eigenvalue-collapse} gives
	\[
	\overline\lambda
	=
	\lambda_*
	=
	\underline\lambda.
	\]
	Every accumulation point of \((\hat\lambda_r)\) therefore equals the same
	value \(\hat\lambda\). Since \((\hat\lambda_r)\) is bounded by
	\eqref{eq:estM}, this proves
	\[
	\hat\lambda_r\to\hat\lambda.
	\]
	
	Finally, let \((\mathbf u_*,\mathbf v_*)\) be any weak-in-\(\mathcal V\) and
	strong-in-\(H\) cluster point of the normalized discrete eigenpairs. Since
	\(\hat\lambda_r\to\hat\lambda\), the preceding passage to the limit applies
	with \(\lambda_*=\hat\lambda\). Hence
	\[
	\mathbf u_*,\mathbf v_*
	\in
	\mathcal S^o,
	\]
	they solve
	\[
	(\mathcal L-\hat\lambda\mathcal G)\mathbf u_*=0,
	\qquad
	(\mathcal L-\hat\lambda\mathcal G)^*\mathbf v_*=0,
	\]
	and satisfy
	\[
	\mathcal R(\mathbf u_*,\mathbf v_*)=\hat\lambda.
	\]
	This proves the theorem.
\end{proof}


\begin{thebibliography}{99}
\bibitem[BerestNienbVardan]{BerestNienbVardan} H. Berestycki, L. Nirenberg and S. R. S. Varadhan, The principal eigenvalue and maximum principle for second-order elliptic operators in general domains, Comm. Pure Appl. Math. 47 (1994), 47--92.
\bibitem[Donsker]{Donsker} M. D. Donsker and S. R. S. Varadhan, On a variational formula for the principal eigenvalue for operators with maximum principle, Proc. Natl. Acad. Sci. USA 72 (1975), 780--783.
\bibitem[Ern]{Ern} A. Ern and J.-L. Guermond, \textit{Theory and Practice of Finite Elements}, Applied Mathematical Sciences, vol. 159, Springer, New York, 2004.
\bibitem[ProtterWeinberger]{ProtterWeinberger} M. H. Protter and H. F. Weinberger, \textit{Maximum Principles in Differential Equations}, Springer, New York, 1984.
\bibitem[Sion1958]{Sion1958} M. Sion, On general minimax theorems, Pacific J. Math. 8 (1958), 171--176.
\bibitem[ciaret]{ciaret} P. G. Ciarlet, \textit{The Finite Element Method for Elliptic Problems}, Studies in Mathematics and its Applications, vol. 4, North-Holland, Amsterdam, 1978.
\bibitem[ciaretRav]{ciaretRav} P. G. Ciarlet and P. A. Raviart, Maximum principle and uniform convergence for the finite element method, Comput. Methods Appl. Mech. Engrg. 2 (1973), 17--31.
\end{thebibliography}
\end{document}
